\section{Executive Summary and Governance Framework}
\label{sec:executive_summary}

\begin{tcolorbox}[colback=blue!5!white,colframe=blue!60!black,title=\textbf{Governing Supervisor Directive \& Active Phase},top=2pt,bottom=2pt,left=4pt,right=4pt]
\small
\emph{``We need to have understood everything related to the first model before we increase complexity.''}
\begin{center}
  \textbf{\texttt{MODE1\_GATE6B\_ACTIVE\_EVALUATION\_AND\_CONTINUATION}}
\end{center}
\end{tcolorbox}

Following the supervisor decisions of 17 September 2026, this report provides the comprehensive Mode-I convergence and energetic audit package for the 08 October 2026 supervisory review. The primary objective is to complete the physical, numerical, and energetic qualification of the Mode-I phase-field benchmark before advancing to state transfer, visualization integration, or multi-crack geometries.

\subsection{Status of governing milestone gates}
\begin{enumerate}[label=\textbf{\arabic*.},leftmargin=6mm,itemsep=2pt]
  \item \textbf{Gate 6A -- Mechanical Implementation \& Defect Closure:} \textbf{\texttt{RESOLVED\_AND\_CLOSED}}. The early 71,320-element initial stiffness defect was proven to arise from an Abaqus keyword preprocessing line limit ($\le 16$ nodes per line in \texttt{*NSET} without \texttt{GENERATE}), silently omitting 134 of 150 boundary nodes. Wrapping the data lines restored all constraints and achieved reference-consistent structural stiffness recovery (within $0.09\%$, $K_0 = 137.820804\,\mathrm{kN/mm}$ in full fracture Job \texttt{1404933}, $\Delta K_0 = -0.09\%$, $\Delta F_{\max} = -1.64\%$).
  \item \textbf{Gate 5 -- Native Remeshing Reproduction (71,320 vs $\approx 13{,}941$):} \textbf{\nolinkurl{SUPERVISOR_ACCEPTED_REPRODUCTION_LIMITATION_CLOSED}}. The supervisor accepted that publication details do not expose enough information to reproduce every reported element count exactly. Parametric sensitivity trends are preserved and strictly disambiguated between offline preliminary sizing ($1.0\% \to 71{,}320$; $2.0\% \to 17{,}687$; $3.0\% \to 8{,}120$; $5.0\% \to 4{,}356$) and the audited production batch meshes ($1.0\% \to 71{,}320$; $2.0\% \to 15{,}396$; $3.0\% \to 7{,}633$; $5.0\% \to 4{,}194$ finite elements).
  \item \textbf{Gate 6B -- Mode-I Energetic and Multifaceted Convergence:} \textbf{\texttt{ACTIVE -- HIGHEST PRIORITY}}. Decoupled evaluation across structural stiffness $K_0$, peak force $F_{\max}$, peak displacement $u_{\mathrm{peak}}$, trapezoidal boundary work $W_{\mathrm{trap}}$, UEL energy terms ($E_{\mathrm{elas}}, E_{\mathrm{frac}}$), crack-path symmetry, temporal discretization ($T_1$--$T_3$ and adaptive $2\times$), spatial discretization ($S_1$--$S_4$), and length scale ($l_0$).
  \begin{itemize}[leftmargin=4mm,itemsep=1pt]
    \item \textbf{UEL Energy Output:} Formally promoted to \textbf{\texttt{UEL\_ENERGY\_OUTPUT\_QUALIFIED\_MECHANICALLY\_NONINVASIVE}} via full 7,000-increment point-by-point mechanical parity audit ($\Delta K_0 = 0.000000\%$, $\Delta F_{\max} = 0.000000\%$).
    \item \textbf{Temporal Convergence:} Pre-peak response and peak characteristics are \textbf{\texttt{TEMPORALLY\_STABLE}} ($\Delta K_0 = +0.000302\%$, $\Delta F_{\max} = -0.0229\%$, \textbf{\texttt{QUALIFIED\_OVER\_PREPEAK\_INTERVAL\_ONLY}}), while post-fracture wake exhibits increment-scale normalization sensitivity (\textbf{\texttt{TEMPORALLY\_SENSITIVE\_POSTPEAK}}).
    \item \textbf{Active Solver Solves on \texttt{/scratch9/}:} Actively tracking spatial candidate Job \texttt{1410179} (58k FE), convergence-control diagnostic Job \texttt{1410180} ($C_n=0.50$), and errorTarget fracture batch Jobs \texttt{1410357}--\texttt{1410359} (ET2, ET3, ET5).
  \end{itemize}
  \item \textbf{Global Energy Identity:} \textbf{\texttt{GLOBAL\_ENERGY\_IDENTITY --- NOT\_YET\_CLOSED}}. The staggered solution scheme evaluates history $\mathcal{H}_n$ at step $n$ and damage $d_{n+1}$ at step $n+1$, disproving the existence of a joint discrete scalar potential ($\partial^2 \Pi / \partial \mathbf{u} \partial d \ne \partial^2 \Pi / \partial d \partial \mathbf{u}$).
  \item \textbf{Active Scope Holds:} Mode-II, Gate~6C (State Transfer), Gate~7 (Visualization integration), and multi-rank MPI remain on \textbf{strict hold} pending supervisor acceptance of Gate~6B closure.
\end{enumerate}

\subsection{Key scientific findings presented in this package}
\begin{itemize}[leftmargin=5mm,itemsep=2pt]
  \item \textbf{Energy Source Mechanical Parity:} Formally qualified as \textbf{\texttt{UEL\_ENERGY\_OUTPUT\_QUALIFIED\_MECHANICALLY\_NONINVASIVE}}. Global mechanical force trajectories match bitwise ($|\Delta F| = 0.000000\,\mathrm{kN}$, $\Delta K_0 = 0.000000\%$) across all $7{,}000$ increments ($21{,}120$ Newton iterations, 0 cutbacks) against uninstrumented baselines.
  \item \textbf{Single-IP Deduplication Rule:} Companion visualizer UMAT elements replicate scalar whole-element energies into all 4 Gauss integration points. Summing across all IPs creates an artificial $4\times$ overcounting error. Enforcing extraction at Integration Point 1 (\texttt{IP1}) recovers the once-per-underlying-finite-element global energy sum under the 2D implicit-unit-thickness convention.
  \item \textbf{Structural Stiffness Invariance:} Initial stiffness $K_0$ is exceptionally stable across spatial meshes ($S_1$--$S_4$: $137.95$ to $137.84\,\mathrm{kN/mm}$, spread $<0.09\%$), native adaptive meshes ($A_1$--$A_4$: spread $<0.09\%$), and time steps ($\Delta K_0 = +0.000302\%$), confirming kinematic, geometric, and elastic consistency.
  \item \textbf{Peak Force Mesh Sensitivity:} Peak reaction force drops monotonically with spatial mesh refinement from $0.7578\,\mathrm{kN}$ ($S_1$, $h/l_0 = 0.400$) to $0.7312\,\mathrm{kN}$ ($S_4$, $h/l_0 = 0.167$), extending to $0.7255\,\mathrm{kN}$ on the historical fine mesh ($h = 1.0\,\mu\mathrm{m}$, $-4.26\%$ total variation), classified as \textbf{\texttt{MESH-SENSITIVE}}.
  \item \textbf{Crack-Path Symmetry and Profile Convergence:} The damage-weighted crack-centroid trajectory deviates from the symmetry plane $y = 0.500\,\mathrm{mm}$ by $\le 3.10\,\mu\mathrm{m}$ across all fixed and adaptive meshes, confirming pure Mode-I symmetry within discrete element bounds.
  \item \textbf{Completed 3-Point Length-Scale Study ($l_0$):} Evaluated across $l_0 \in \{7.5, 11.25, 15.0\}\,\mu\mathrm{m}$ on fixed $S_3$ mesh (Jobs \texttt{1406017.mmaster02}, \texttt{1406895.mmaster02}, \texttt{1406896.mmaster02}). $K_0$ is stable ($137.86 \to 137.68\,\mathrm{kN/mm}$, $0.13\%$ spread), while $F_{\max}$ ($0.7322 \to 0.6895\,\mathrm{kN}$) and $u_{\mathrm{peak}}$ ($5.633 \to 5.579\,\mu\mathrm{m}$) are length-scale sensitive.
\end{itemize}
